% BHU PYQ -- Manually transcribed & error-corrected
% Paper: BCH-216 : Basic Statistics
% Exam: B.Com. (Hons.) Semester III Examination, 2022-23

\documentclass[11pt,a4paper]{article}
\usepackage[a4paper,top=2.5cm,bottom=2.5cm,left=2.8cm,right=2.8cm,headheight=14pt]{geometry}
\usepackage{amsmath,amssymb}
\usepackage[T1]{fontenc}
\usepackage[utf8]{inputenc}
\usepackage{lmodern,microtype,fancyhdr,enumitem,hyperref,lastpage,parskip,booktabs}

\hypersetup{
    colorlinks=true,
    urlcolor=black,
    linkcolor=black,
    pdftitle={BCH-216: Basic Statistics -- BHU B.Com. (Hons.) Sem III 2022-23}
}

\pagestyle{fancy}
\fancyhf{}
\renewcommand{\headrulewidth}{0.4pt}
\renewcommand{\footrulewidth}{0.4pt}
\fancyhead[L]{\small   \textit{Banaras Hindu University}}
\fancyhead[R]{\small   \textit{BCH-216: Basic Statistics}}
\fancyfoot[C]{\small Page  \thepage\ of \pageref{LastPage}}


\newlist{parts}{enumerate}{2}
\setlist[parts,1]{label=(\alph*),leftmargin=*,itemsep=5pt,topsep=3pt}
\setlist[parts,2]{label=(\roman*),leftmargin=*,itemsep=3pt,topsep=2pt}

\newcommand{\pts}[1]{\hfill   \textit{[#1]}}


\begin{document}


\begin{center}
  {\large\bfseries BANARAS HINDU UNIVERSITY}\\[2pt]
  {\normalsize B.Com. (Hons.) --- Semester III Examination, 2022-23}\\[6pt]
  \rule{0.8\linewidth}{0.5pt}\\[6pt]
  {\large\bfseries Paper No. BCH-216: Basic Statistics}\\[4pt]
  {\normalsize Time: 3 Hours\hfill Maximum Marks: 70}
\end{center}

\rule{\linewidth}{0.4pt}

\smallskip



\noindent   \textit{Instructions: Answer all questions. All questions carry equal marks (14 marks each).}

\bigskip




\noindent   \textbf{Question 1.} \\
``The successful businessman is the one whose estimates most closely approaches accuracy.'' Explain and discuss the limitations of Statistics.\pts{14}

\centerline{   \textbf{OR}}

Draw a cumulative frequency curve of the distribution given below and read the values of the median and two quartiles:

\begin{center}
  
\begin{tabular}{ccccccc}
     oprule
   \textbf{Marks} & 0--10 & 10--20 & 20--30 & 30--40 & 40--50 & 50--60 \\
    \midrule
   \textbf{Frequency} & 4 & 11 & 15 & 24 & 10 & 6 \\
    ottomrule
  \end{tabular}
\end{center}\pts{14}

\medskip\rule{\linewidth}{0.2pt}




\noindent   \textbf{Question 2.} \\
Explain the relationship among arithmetic mean, geometric mean and harmonic mean.\pts{14}

\centerline{   \textbf{OR}}


\begin{parts}
  \item Gains and losses (in Rs.) of 50 speculators are given below:
  
\begin{itemize}
    \item 4 speculators gain between Rs. 300 and Rs. 400
    \item 6 speculators gain between Rs. 200 and Rs. 300
    \item 8 speculators gain between Rs. 100 and Rs. 200
    \item 12 speculators gain between Rs. 0 and Rs. 100
    \item 10 speculators lose between Rs. 0 and Rs. 100
    \item 7 speculators lose between Rs. 100 and Rs. 200
    \item 3 speculators lose between Rs. 200 and Rs. 300
  \end{itemize}
  Find out the average gain of gainers, average loss of losers and average net gain or loss of the whole group.\pts{7}
  
  \item Compute the mean and mode marks of all the students of 50 colleges in a District given below:
  
\begin{center}
    \small
    
\begin{tabular}{ccc}
       oprule
   \textbf{Marks obtained} &   \textbf{No. of colleges} &   \textbf{Average number of students in a college} \\
      \midrule
      35 and above & 7 & 200 \\
      30--35 & 10 & 250 \\
      25--30 & 15 & 300 \\
      20--25 & 9 & 200 \\
      15--20 & 5 & 150 \\
      Less than 15 & 4 & 100 \\
      ottomrule
    \end{tabular}
  \end{center}\pts{7}
\end{parts}

\medskip\rule{\linewidth}{0.2pt}




\noindent   \textbf{Question 3.} \\
The following data are related with the mileage of two types of tyres:

\begin{center}
  
\begin{tabular}{ccc}
     oprule
   \textbf{Life (in km)} &   \textbf{No. of tyres: Type A} &   \textbf{No. of tyres: Type B} \\
    \midrule
    20000--22000 & 230 & 200 \\
    22000--24000 & 270 & 325 \\
    24000--26000 & 450 & 470 \\
    26000--28000 & 375 & 300 \\
    28000--30000 & 125 & 155 \\
    ottomrule
  \end{tabular}
\end{center}

\begin{parts}
  \item Which of the two types give a higher average life?\pts{7}
  \item If the price of these two types are same, which one will you prefer and why?\pts{7}
\end{parts}

\centerline{   \textbf{OR}}

Calculate Karl Pearson's coefficient of skewness from the following data:

\begin{center}
  
\begin{tabular}{lr}
     oprule
   \textbf{Marks} &   \textbf{No. of students} \\
    \midrule
    Less than 20 & 8 \\
    Less than 30 & 38 \\
    Less than 40 & 53 \\
    40--50 & 2 \\
    50--60 & 8 \\
    60--70 & 30 \\
    More than 70 & 17 \\
    More than 80 & 4 \\
    ottomrule
  \end{tabular}
\end{center}\pts{14}

\medskip\rule{\linewidth}{0.2pt}




\noindent   \textbf{Question 4.} \\
For 50 paired observations of $X$ and $Y$, the following data were obtained:
\[ ar{X} = 10, \quad \sigma_x = 3, \quad ar{Y} = 6, \quad \sigma_y = 2, \quad r_{xy} = 0.3 \]
But on subsequent verification it was found that in a pair one value of $X (= 10)$ and one value of $Y (= 6)$ were inaccurate and hence needed out. With the remaining 49 pairs of values, how is the original value of correlation coefficient affected?\pts{14}

\centerline{   \textbf{OR}}

Distinguish clearly among simple, partial and multiple correlation and point out their usefulness in statistical analysis.\pts{14}

\medskip\rule{\linewidth}{0.2pt}




\noindent   \textbf{Question 5.} \\
Explain regression coefficients. How are the regression coefficients of two regression lines related with the correlation coefficient?\pts{14}

\centerline{   \textbf{OR}}

From the following data:

\begin{center}
  
\begin{tabular}{cc}
     oprule
   \textbf{Age of husband (Year)} &   \textbf{Age of wife (Year)} \\
    \midrule
    18 & 17 \\
    19 & 17 \\
    20 & 18 \\
    21 & 18 \\
    22 & 18 \\
    23 & 19 \\
    24 & 19 \\
    25 & 20 \\
    26 & 21 \\
    27 & 21 \\
    ottomrule
  \end{tabular}
\end{center}
find:

\begin{parts}
  \item the regression equation of age of husband on the age of wife;\pts{4}
  \item the regression equation of age of wife on the age of husband;\pts{4}
  \item total variation in the age of husband;\pts{3}
  \item the magnitude of variation in age of the husband explained by the regression equation.\pts{3}
\end{parts}

\vfill
\rule{\linewidth}{0.4pt}

\begin{center}\small
  \href{https://bhu-exam.vercel.app/}{Click here to visit our website}\\[4pt]
  \href{https://me-aryan.vercel.app/}{Click here to visit our contributor}\\[4pt]
  \href{https://quashan-me.vercel.app/}{Click here to visit our contributor}
\end{center}

\end{document}
